\chapter{Circuitos em série e em paralelo}\label{ch:g5:series-parallel-circuits}

Quando uma lâmpada queimava, o cordão de festa inteiro apagava --- a
maldição da volta única. E, no entanto, quando uma lâmpada queima na
sua cozinha, o resto da casa continua aceso. Alguém ligou a sua casa
de um jeito mais esperto. Hoje aprendemos esse jeito: o jeito dos
\omterm{def:g5:series-parallel-circuits:parallel}{ramos}.

\section{Dois jeitos de ligar}

Você já conhece o primeiro: \emph{\omterm{def:g4:series-circuits:series}{em série}}, com todos os jogadores
enfiados numa volta só, contas num mesmo colar. Aqui vem o segundo.

\begin{definition}[Em paralelo]\label{def:g5:series-parallel-circuits:parallel}
Os jogadores estão \emph{em paralelo}\index{circuito em paralelo}
quando o circuito se divide em \emph{ramos}\index{ramo} separados ---
um jogador por ramo --- que se separam e depois se juntam de novo. Os
pontos em que as estradas se separam ou se juntam se chamam
\emph{nós}\index{nó}. Cada ramo é uma estrada completa de um \omterm{def:g3:first-electric-circuit:battery}{terminal}
da \omterm{def:g3:first-electric-circuit:battery}{pilha} até o outro.
\end{definition}

\begin{omfigure}
\begin{tikzpicture}[scale=0.8]
  \draw (0,0) to[battery1, l=pilha] (0,3.2)
    to[short] (2.2,3.2);
  \draw (2.2,3.2) to[short] (2.2,2.4)
    to[lamp, l=lâmpada 1] (5.4,2.4) to[short] (5.4,3.2);
  \draw (2.2,3.2) to[short] (2.2,4.0)
    to[lamp, l=lâmpada 2] (5.4,4.0) to[short] (5.4,3.2);
  \draw (5.4,3.2) to[short] (7.0,3.2) to[short] (7.0,0)
    to[short] (0,0);
  \fill (2.2,3.2) circle (2pt);
  \fill (5.4,3.2) circle (2pt);
  \node[above, font=\small] at (3.6,4.95) {dois ramos};
\end{tikzpicture}

{\small Duas lâmpadas \omterm{def:g5:series-parallel-circuits:parallel}{em paralelo}: a estrada se divide num \omterm{def:g5:series-parallel-circuits:parallel}{nó}, cada
lâmpada viaja no seu próprio \omterm{def:g5:series-parallel-circuits:parallel}{ramo}, e os \omterm{def:g5:series-parallel-circuits:parallel}{ramos} se juntam de novo antes
de voltar para a \omterm{def:g3:first-electric-circuit:battery}{pilha}.}
\end{omfigure}

\begin{method}[Montar o circuito]\label{met:g5:series-parallel-circuits:build}
Com uma \omterm{def:g3:first-electric-circuit:battery}{pilha}, duas lâmpadas e fios:
\begin{enumerate}
  \item ligue a primeira lâmpada à \omterm{def:g3:first-electric-circuit:battery}{pilha}, sozinha, como uma volta
        simples --- \omterm{def:g5:series-parallel-circuits:parallel}{ramo} um;
  \item agora ligue a segunda lâmpada \emph{diretamente sobre} a
        primeira: um fio de cada lado da lâmpada 1 até o lado
        correspondente da lâmpada 2 --- \omterm{def:g5:series-parallel-circuits:parallel}{ramo} dois;
  \item percorra cada \omterm{def:g5:series-parallel-circuits:parallel}{ramo} com o dedo: saindo do $+$, toda estrada
        deve passar por exatamente uma lâmpada antes de chegar ao
        $-$.
\end{enumerate}
As duas lâmpadas acendem --- e as duas brilham tão forte quanto uma
lâmpada sozinha brilharia.
\end{method}

\section{As regras dos ramos}

\begin{proposition}[Regras do circuito em paralelo]\label{prop:g5:series-parallel-circuits:rules}
Com cada jogador no seu próprio \omterm{def:g5:series-parallel-circuits:parallel}{ramo}:
\begin{enumerate}
  \item \emph{independência}: uma falha num \omterm{def:g5:series-parallel-circuits:parallel}{ramo} para só aquele \omterm{def:g5:series-parallel-circuits:parallel}{ramo}
        --- desrosqueie uma lâmpada e a outra continua acesa;
  \item \emph{brilho total}: cada lâmpada brilha como se estivesse
        sozinha --- não há repartição do empurrão da \omterm{def:g3:first-electric-circuit:battery}{pilha} entre os
        \omterm{def:g5:series-parallel-circuits:parallel}{ramos};
  \item \emph{postos de comando}: um \omterm{ex:g3:first-electric-circuit:switch}{interruptor} dentro de um \omterm{def:g5:series-parallel-circuits:parallel}{ramo}
        manda só naquele \omterm{def:g5:series-parallel-circuits:parallel}{ramo}; um \omterm{ex:g3:first-electric-circuit:switch}{interruptor} na estrada comum, antes
        da divisão, manda em todo mundo. (O preço: a \omterm{def:g3:first-electric-circuit:battery}{pilha} alimenta
        todos os \omterm{def:g5:series-parallel-circuits:parallel}{ramos} ao mesmo tempo e, por isso, se descarrega mais
        depressa.)
\end{enumerate}
\end{proposition}

\begin{example}[O teste da desrosqueada, de novo]\label{ex:g5:series-parallel-circuits:test}
Faça com a dupla \omterm{def:g5:series-parallel-circuits:parallel}{em paralelo} exatamente o que um dia matou a dupla \omterm{def:g4:series-circuits:series}{em
série}: desrosqueie a lâmpada 2. A lâmpada 1 nem pisca --- o \omterm{def:g5:series-parallel-circuits:parallel}{ramo}
dela, a estrada completa dela, continua intacto. Rosqueie de volta e
desrosqueie a lâmpada 1: a lâmpada 2 continua acesa, tranquila. A
maldição da volta única foi quebrada: \omterm{def:g5:series-parallel-circuits:parallel}{em paralelo}, cada jogador cuida
do próprio \omterm{def:g5:series-parallel-circuits:parallel}{ramo}.
\end{example}

\begin{omfigure}
\begin{tikzpicture}[scale=0.8]
  \draw (0,0) to[battery1] (0,3.2)
    to[cspst, l=interruptor geral] (2.6,3.2);
  \draw (2.6,3.2) to[short] (2.6,2.4)
    to[lamp, l=lâmpada 1] (5.8,2.4) to[short] (5.8,3.2);
  \draw (2.6,3.2) to[short] (2.6,4.0)
    to[spst, l=interruptor do ramo] (4.2,4.0)
    to[lamp] (5.8,4.0) to[short] (5.8,3.2);
  \draw (5.8,3.2) to[short] (7.4,3.2) to[short] (7.4,0)
    to[short] (0,0);
  \fill (2.6,3.2) circle (2pt);
  \fill (5.8,3.2) circle (2pt);
\end{tikzpicture}

{\small Postos de comando: o \omterm{ex:g3:first-electric-circuit:switch}{interruptor} geral, antes da divisão,
manda nos dois \omterm{def:g5:series-parallel-circuits:parallel}{ramos}; o \omterm{ex:g3:first-electric-circuit:switch}{interruptor} do \omterm{def:g5:series-parallel-circuits:parallel}{ramo} manda só na lâmpada dele
--- aqui ele está aberto, então a lâmpada 1 brilha sozinha.}
\end{omfigure}

\begin{example}[A sua casa é um circuito em paralelo]\label{ex:g5:series-parallel-circuits:house}
Cada luminária, tomada e aparelho da sua casa viaja no próprio \omterm{def:g5:series-parallel-circuits:parallel}{ramo}
de um grande \omterm{def:g5:series-parallel-circuits:parallel}{circuito em paralelo}. É por isso que a lâmpada da
cozinha pode queimar enquanto a sala continua com \omterm{def:g1:light-and-shadows:source}{luz} para ler; por
isso que o \omterm{ex:g3:first-electric-circuit:switch}{interruptor} de cada cômodo manda só na \omterm{def:g1:light-and-shadows:source}{luz} daquele cômodo;
e por isso que um único \omterm{ex:g3:first-electric-circuit:switch}{interruptor} geral --- no quadro de \omterm{def:g2:pushes-and-pulls:force}{força}, ao
lado do medidor --- corta a casa inteira de uma vez para consertos.
Os cordões de luzinhas modernos entraram para o clube: as lâmpadas
deles são ligadas de modo que uma lâmpada queimada não apaga mais o
cordão.
\end{example}

\begin{remark}[Série e paralelo, lado a lado]\label{rem:g5:series-parallel-circuits:compare}
Uma falha: \omterm{def:g4:series-circuits:series}{em série} --- tudo apagado; \omterm{def:g5:series-parallel-circuits:parallel}{em paralelo} --- só o \omterm{def:g5:series-parallel-circuits:parallel}{ramo} dela
apagado. Mais lâmpadas: \omterm{def:g4:series-circuits:series}{em série} --- todas mais fracas; \omterm{def:g5:series-parallel-circuits:parallel}{em paralelo}
--- todas no brilho total, mas a \omterm{def:g3:first-electric-circuit:battery}{pilha} se descarrega mais depressa.
Um \omterm{ex:g3:first-electric-circuit:switch}{interruptor} para tudo: \omterm{def:g4:series-circuits:series}{em série} --- em qualquer ponto; \omterm{def:g5:series-parallel-circuits:parallel}{em paralelo}
--- só antes da divisão. Nenhum dos dois jeitos é ``melhor'': a
lanterna é feliz \omterm{def:g4:series-circuits:series}{em série} (uma lâmpada, um \omterm{ex:g3:first-electric-circuit:switch}{interruptor}), a casa é
obrigatoriamente \omterm{def:g5:series-parallel-circuits:parallel}{em paralelo}. A primeira pergunta de um eletricista
diante de qualquer circuito é simplesmente: \emph{onde a estrada se
divide?}
\end{remark}

\section{Ler circuitos}

\begin{method}[Percorrer as estradas]\label{met:g5:series-parallel-circuits:follow}
Para prever o destino de qualquer lâmpada num desenho de circuito:
\begin{enumerate}
  \item ponha o dedo no $+$ da \omterm{def:g3:first-electric-circuit:battery}{pilha} e percorra todas as estradas
        possíveis até o $-$, listando os jogadores por que cada
        estrada passa;
  \item uma lâmpada acende se pelo menos uma estrada sem
        interrupções passar por ela;
  \item jogadores que dividem todas as estradas estão \omterm{def:g4:series-circuits:series}{em série} entre
        si; jogadores em divisões diferentes estão \omterm{def:g5:series-parallel-circuits:parallel}{em paralelo}.
\end{enumerate}
Percorra as estradas antes de prever --- chutar só pelo formato do
desenho é o jeito como os circuitos enganam as pessoas.
\end{method}

\section{Exercícios}

\begin{exercise}[$\star$]\label{exo:g5:series-parallel-circuits:1}
O que é um \omterm{def:g5:series-parallel-circuits:parallel}{ramo}? E um \omterm{def:g5:series-parallel-circuits:parallel}{nó}? Num \omterm{def:g5:series-parallel-circuits:parallel}{circuito em paralelo} com duas lâmpadas,
quantos \omterm{def:g5:series-parallel-circuits:parallel}{ramos} carregam uma lâmpada?
\end{exercise}

\begin{exercise}[$\star$]\label{exo:g5:series-parallel-circuits:2}
Enuncie as três regras do \omterm{def:g5:series-parallel-circuits:parallel}{circuito em paralelo}.
\end{exercise}

\begin{exercise}[$\star$]\label{exo:g5:series-parallel-circuits:3}
Duas lâmpadas \omterm{def:g5:series-parallel-circuits:parallel}{em paralelo}: a lâmpada 1 é desrosqueada. O que a
lâmpada 2 faz? E na dupla \omterm{def:g4:series-circuits:series}{em série} --- o que ela fazia lá, e por que
a diferença?
\end{exercise}

\begin{exercise}[$\star$]\label{exo:g5:series-parallel-circuits:4}
Duas lâmpadas iguais: a dupla A está \omterm{def:g4:series-circuits:series}{em série}, a dupla B \omterm{def:g5:series-parallel-circuits:parallel}{em paralelo},
com \omterm{def:g3:first-electric-circuit:battery}{pilhas} iguais. Qual dupla brilha mais forte por lâmpada? Qual
descarrega a \omterm{def:g3:first-electric-circuit:battery}{pilha} mais depressa?
\end{exercise}

\begin{exercise}[$\star$]\label{exo:g5:series-parallel-circuits:5}
Na figura com os dois \omterm{ex:g3:first-electric-circuit:switch}{interruptores}, o que cada lâmpada faz quando:
os dois \omterm{ex:g3:first-electric-circuit:switch}{interruptores} estão fechados? \ o geral está aberto? \ o
geral está fechado e o do \omterm{def:g5:series-parallel-circuits:parallel}{ramo}, aberto?
\end{exercise}

\begin{exercise}[$\star$]\label{exo:g5:series-parallel-circuits:6}
Por que uma casa é ligada \omterm{def:g5:series-parallel-circuits:parallel}{em paralelo} e não \omterm{def:g4:series-circuits:series}{em série}? Dê dois
confortos do dia a dia que dependem disso.
\end{exercise}

\begin{exercise}[$\star$]\label{exo:g5:series-parallel-circuits:7}
Em que ponto do circuito de uma casa um \omterm{ex:g3:first-electric-circuit:switch}{interruptor} manda em tudo de
uma vez? Qual é o nome e o uso desse \omterm{ex:g3:first-electric-circuit:switch}{interruptor} no dia a dia?
\end{exercise}

\begin{exercise}[$\star$]\label{exo:g5:series-parallel-circuits:8}
Um abajur e a tomada dele: eles estão \omterm{def:g4:series-circuits:series}{em série} ou \omterm{def:g5:series-parallel-circuits:parallel}{em paralelo} com a
\omterm{def:g1:light-and-shadows:source}{luz} do teto? Como você sabe disso pelo dia a dia?
\end{exercise}

\begin{exercise}[$\star\star$]\label{exo:g5:series-parallel-circuits:9}
Desenhe (com símbolos) um circuito para um corredor: uma \omterm{def:g3:first-electric-circuit:battery}{pilha}, uma
lâmpada e este conforto --- a lâmpada precisa acender e apagar por um
\omterm{ex:g3:first-electric-circuit:switch}{interruptor} em \emph{qualquer} uma das duas pontas do corredor\dots\
Tente de verdade; se os seus projetos insistirem em falhar, explique
o que é difícil. (O verdadeiro comando de dois pontos precisa de um
\omterm{ex:g3:first-electric-circuit:switch}{interruptor} mais esperto que os nossos --- a tentativa ensina mais
que a resposta.)
\end{exercise}

\begin{exercise}[$\star\star$]\label{exo:g5:series-parallel-circuits:10}
Um cordão de luzinhas modernas continua aceso quando uma lâmpada
queima --- mas a caixa avisa: ``Não use o cordão com mais de três
lâmpadas retiradas.'' Usando as regras dos \omterm{def:g5:series-parallel-circuits:parallel}{ramos}, explique o conforto
e, mais difícil, um motivo sensato para o aviso. (Pense no que cada
\omterm{def:g5:series-parallel-circuits:parallel}{ramo} restante recebe da \omterm{def:g3:first-electric-circuit:battery}{pilha}, e no esforço da própria \omterm{def:g3:first-electric-circuit:battery}{pilha}.)
\end{exercise}

\begin{exercise}[$\star\star$]\label{exo:g5:series-parallel-circuits:11}
Na figura dos dois \omterm{ex:g3:first-electric-circuit:switch}{interruptores}, um eletricista quer acrescentar
mais uma lâmpada que acenda \emph{somente} quando o \omterm{ex:g3:first-electric-circuit:switch}{interruptor} do
\omterm{def:g5:series-parallel-circuits:parallel}{ramo} estiver fechado, e fique apagada nos outros casos --- mesmo com
o \omterm{ex:g3:first-electric-circuit:switch}{interruptor} geral fechado. Em qual trecho de estrada a lâmpada nova
precisa ser ligada, e ela fica então \omterm{def:g4:series-circuits:series}{em série} ou \omterm{def:g5:series-parallel-circuits:parallel}{em paralelo} com a
lâmpada 2?
\end{exercise}
